\section{Artifact identity and reproduction}
\label{app:artifact}
The repository is \url{https://github.com/KSwordDEV/KSword}. The source revisions
used in this preprint are:
\begin{center}\footnotesize
\begin{tabular}{@{}ll@{}}
Evidence snapshot & \code{d0d15798f1d92cbdb66244370e6c0c2d4487feb9} \\
Current implementation & \code{bd703d84a0f45f7567c8d05113578ae9e4adf1bd} \\
Historical optimized & \code{b9bdda8358fc7278e79053f2cc4dbf2056a900b9} \\
\end{tabular}
\end{center}
These full identifiers identify source revisions, not merely a branch name.
\code{provenance.json} and the historical build provenance additionally record
source-file byte hashes and measured binary identities. The current driver and
controller SHA256 values, respectively, are:
\begin{center}\footnotesize
\code{3b44132f82febce3fcdf059b697d938479b9f791a628c00b55bc2cf397d5d500}\\[3pt]
\code{4dd0d8dabb3019503e15d70692759ff0502cac46635d53f266a577caad5972ae}
\end{center}
The historical optimized driver SHA256 is:
\begin{center}\footnotesize
\code{8850e28f7f43cf12f95116c862e0aa9fb4b609d60446ee1d46ddc1b900ab7a04}.
\end{center}

\subsection{Evidence locations}
All paths below are relative to the repository root. The current dataset is
\path{docs/next/paper-data/20260916-followup/}. Its raw index binds 268
non-derived records; the dataset includes environment identities, per-run data,
analysis summaries, measured scripts, and validation results. Its key groups are:
\begin{itemize}
\item \path{attribution/}, \path{attribution-runs.csv}, and
      \path{analysis-followup.json}: same-binary performance runs and summaries.
\item \path{latency/} and \path{latency-runs.csv}: every TCP request and each
      control call, identities, overlap gaps, and internal timing.
\item \path{lifecycle/}: individual map/remove/fault/rejection records, serial
      observations, event traces, and \path{derived/smp-summary.json}.
\item \path{application/}: all four EPT application attempts, the three
      direct-write records, setup failures, exact measured scripts, and the
      strict HTTP and policy-comparison summaries.
\item \path{hyperv-*} records: fixture attempts, normal-init serial, root
      capability refusal, and identity snapshots around admission.
\item \path{outer-environment.json}, \path{outer-cpu-identity.json},
      \path{outer-deviceguard.json}, and \path{provenance.json}: platform,
      HVCI, source, and binary identities.
\end{itemize}

The historical cohort is \path{docs/next/paper-data/20260915-4x2/}.
\path{derived/nested-improvements.csv} and
\path{after/derived/windows-overhead.csv} supply the historical performance
tables. \path{smp/derived/smp-summary.json} supplies its separate idle/load/fault
results. \path{after/derived/internal-transition-cpus.csv} retains CPU phase
times. \path{after-nested/counter-interval.json} preserves the counter-window
boundaries. Historical records are not relabeled as current-build tests.

\subsection{Offline analysis}
Python 3.10 or later and the standard library suffice for the existing evidence
analysis. From the repository root, the following PowerShell commands regenerate
the current summaries without running a VM:
\begin{center}\begin{minipage}{\linewidth}
\begin{verbatim}
$d = "docs/next/paper-data/20260916-followup"
python tools/hvm_paper/analyze_followup.py $d
python tools/hvm_paper/analyze_smp.py "$d/lifecycle"
python tools/hvm_paper/analyze_http_page.py "$d/application"
python tools/hvm_paper/analyze_policy_comparison.py "$d/application"
python docs/next/arxiv/generate_tables.py
\end{verbatim}
\end{minipage}\end{center}
Successful analysis means the stated predicates hold for the available records;
it does not erase incomplete trials or prove untested runtime paths. The
preprint's \path{table-provenance.json} records SHA256 values for consumed
summaries and computed headline ratios. TeX inputs for every numeric table are
included with the manuscript, so paper compilation does not require rerunning
experiments or fetching evidence.

\begin{samepage}
The separate redacted evidence/analysis archive contains 1557 files and is
11,122,891 bytes. Its SHA256 is:
\begin{center}\footnotesize
\code{d90ed459a5937000ff0df8d148086d2d603405dab2cd6726121fe62a8fdf6bef}.
\end{center}
\end{samepage}
It supports offline recomputation and includes an integrity manifest; private
identity mappings are excluded. It is distinct from the arXiv TeX source
package and is not a complete anonymous buildable source artifact. The
repository contains the unredacted versioned evidence. No archive upload or
external artifact acceptance is claimed by this draft.

\subsection{Example continuity identities}
In the current application series, Windows~1 boot time is
\code{2026-09-16T01:23:55.335060Z}; VMware VMX PID is 10748, created at
\code{2026-09-16T01:38:09.239494Z}. The TinyCore boot ID is
\begin{center}\small\code{e4b22249-5b20-4048-b466-f308c7e04a85},\end{center}
and the HTTP server is PID 4578 with creation tick 52514. The raw per-run
snapshots, per-response identities, and guest uptimes delimit the supported
continuity intervals. The later lab cleanup reboot does not extend them.

The separate Hyper-V baseline has Windows boot time
\code{2026-09-16T00:53:15.158483Z}, worker PID 2824 with creation time
\code{2026-09-16T01:17:54.101516Z}, and TinyCore boot ID
\begin{center}\small\code{c6c03f47-16bc-4d0f-9dae-14194fa93ea9}.\end{center}
The child VM GUID remains
\code{74eb68a3-c16a-4e98-827a-05d363bdf3f6}. These identities substantiate an
unchanged baseline during a refused operation, not successful interposition.

\section{Historical Windows benchmark suite}
\label{app:historical}
Table~\ref{tab:windows-history} reports the historical Windows-only comparison
in full. Off/pre and off/post each contribute seven measured runs; on contributes
seven. The reference is the combined fourteen off-mode runs. Negative point
estimates with intervals spanning zero do not demonstrate acceleration or zero
cost. Intervals are descriptive within the measured machine and boot.

\begin{table}[htbp]
\centering\small
\caption{Historical Windows overhead. These values belong to the historical
driver/boot and are not pooled with the current same-binary experiment.}
\label{tab:windows-history}
% Generated by generate_tables.py; do not edit by hand.
\begin{tabular}{@{}lrr@{}}
\toprule
Workload & Elapsed overhead & Bootstrap 95\% interval \\
\midrule
CPUID leaf 0 & +804.95\% & +787.73\% to +822.05\% \\
Direct disk read & +3.26\% & -3.49\% to +10.12\% \\
Direct disk write & -5.55\% & -22.38\% to +11.25\% \\
Integer loop & -1.01\% & -5.19\% to +0.58\% \\
Memory copy & -1.23\% & -8.59\% to +13.96\% \\
Pointer chase & -0.69\% & -1.75\% to +2.76\% \\
TCP bulk & -1.89\% & -11.80\% to +18.51\% \\
TCP RTT & +34.82\% & +29.95\% to +37.65\% \\
\bottomrule
\end{tabular}

\end{table}

\newpage
\section{Compatibility and remaining validation matrix}
\label{app:matrix}
The highest-priority missing experiments are state handoff beneath an active
VMM, current-build source-drift and owner-exit stress, a corrected per-response
holder-liveness collector, an independent external application observer, and
matched active-lease performance on additional machines. None of these is
silently inferred from a successful marker-page demonstration.

\begin{center}\small
\captionsetup{hypcap=false}
\captionof{table}{Evidence scope as of this draft. ``Not measured'' is distinct from
an observed refusal. Each successful result is limited to its recorded version.}
\label{tab:matrix}
\begin{tabularx}{\linewidth}{@{}>{\raggedright\arraybackslash}p{3.5cm}>{\raggedright\arraybackslash}p{4.2cm}>{\raggedright\arraybackslash}X@{}}
\toprule
Configuration & Operation & Status \\
\midrule
Hyper-V / Windows 1, VMM absent & Enter and exit KSword residency & Measured; same observed Windows boot/processes \\
KSword / VMware / TinyCore, 4-by-2 & Later startup; page remap, restore, retirement & Current 3/3 cycles; historical 20/20 including load \\
Already-running VMware / TinyCore & Insert beneath whole running stack & Not implemented; ownership gate retained \\
Inner Hyper-V / TinyCore, 4-by-2 & Normal boot; KSword admission & Baseline boots; VMX unavailable and admission refused \\
Inner Hyper-V descendant & EPT control under inserted KSword & Not reached \\
Windows 10 LTSC descendant & Compatibility and page control & Deferred; not measured \\
Other CPU / Windows builds & Portability & Not measured \\
Current 4-by-2 build & Ten-minute or overnight soak & Not newly measured \\
Arbitrary page reuse / DMA / hardware faults & General lifecycle and atomicity & Outside demonstrated contract \\
\bottomrule
\end{tabularx}
\end{center}
